% Spickzettel – 1 Seite A4 – pdfLaTeX (TeXstudio + MiKTeX)
\documentclass[11pt]{extarticle}
\usepackage[T1]{fontenc}
\usepackage[utf8]{inputenc}
\usepackage[ngerman]{babel}
\usepackage{lmodern}
\usepackage[a4paper,margin=12mm]{geometry}
\usepackage{amsmath,amssymb}
\usepackage{siunitx}\sisetup{locale=DE}
\usepackage{booktabs}
\usepackage{multicol}
\usepackage{enumitem}\setlist{nosep,leftmargin=*}
\usepackage{xcolor}
\definecolor{akzent}{RGB}{0,82,147}
\usepackage{titlesec}
\titleformat{\section}{\color{akzent}\bfseries\normalsize}{}{0pt}{}
\titlespacing*{\section}{0pt}{5pt}{2pt}
\pagestyle{empty}
\setlength{\parindent}{0pt}
\setlength{\tabcolsep}{3pt}
\setlength{\columnsep}{8mm}

\begin{document}
{\Large\color{akzent}\textbf{Spickzettel: Voltmeter, Amperemeter \& ADC}}\hfill Präsentation 02.10.2026\par
\vspace{1mm}\hrule\vspace{2mm}

\fcolorbox{akzent}{akzent!6}{\parbox{\dimexpr\textwidth-2\fboxsep-2\fboxrule}{%
\textbf{1.} Voltmeter \emph{parallel \& hochohmig} -- Amperemeter \emph{in Reihe \& niederohmig}\quad
\textbf{2.} ADC = Abtasten $\to$ Quantisieren $\to$ Codieren\quad
\textbf{3.} Im Multimeter: jede Messung = Spannungsmessung + ADC}}

\begin{multicols}{2}
\section*{Messgeräte}
\small\begin{tabular}{@{}lll@{}}
\toprule
 & Voltmeter & Amperemeter\\ \midrule
Anschluss & parallel & in Reihe\\
ideal $R_i$ & $\infty$ & 0\\
real (DMM) & $\approx\SI{10}{\mega\ohm}$ & Shunt (m$\Omega$\dots$\Omega$)\\
Fehler & Belastungsfehler & Bürdenspannung\\
falsch & Reihe: geht nicht & parallel: \textbf{Kurzschluss!}\\
\bottomrule
\end{tabular}\normalsize

\smallskip
$n=\frac{\text{neu}}{\text{alt}}$,\quad $R_V=(n-1)R_i$,\quad $R_N=\dfrac{R_i}{n-1}$\\
Kleines $R$ $\to$ spannungsrichtig, großes $R$ $\to$ stromrichtig

\section*{Sicherheit}
\begin{itemize}
\item Vor jeder Messung: \textbf{Buchse -- Drehschalter -- Bereich}
\item CAT IV Hausanschluss \textbar{} CAT III Verteiler \textbar{} CAT II Steckergeräte \textbar{} CAT I/0 ohne Netz
\item Niedrigste Kategorie (Gerät/Leitung) gilt
\end{itemize}

\section*{Wandlertypen}
\begin{itemize}
\item \textbf{Flash}: sehr schnell, grob, $2^n-1$ Komparatoren $\to$ Oszilloskop
\item \textbf{SAR}: Bit für Bit, $n$ Schritte $\to$ Mikrocontroller
\item \textbf{Dual-Slope}: langsam, genau, störfest $\to$ Multimeter
\item \textbf{Sigma-Delta}: sehr hohe Auflösung $\to$ Audio, Waagen
\end{itemize}

\columnbreak
\section*{ADC-Formeln (Beispiel 10 Bit, 5 V)}
\small\begin{tabular}{@{}lll@{}}
\toprule
Stufen & $2^n$ & 1024\\
LSB & $U_\text{ref}/2^n$ & \SI{4.88}{\milli\volt}\\
Code & $\lfloor U_\text{ein}2^n/U_\text{ref}\rfloor$ & \SI{2.5}{\volt} $\to$ 512\\
Spannung & $\text{Code}\cdot U_\text{ref}/2^n$ & 300 $\to$ \SI{1.46}{\volt}\\
Q-Fehler & $\pm\frac12$\,LSB & $\pm\SI{2.44}{\milli\volt}$\\
Nyquist & $f_A>2f_\text{max}$ & CD \SI{44.1}{\kilo\hertz}\\
Alias & $|f_A-f_\text{sig}|$ & $15-10=\SI{5}{\kilo\hertz}$\\
\bottomrule
\end{tabular}\normalsize\\[1mm]
8\,Bit \SI{19.5}{\milli\volt} \textbar{} 10\,Bit \SI{4.88}{\milli\volt} \textbar{} 12\,Bit \SI{1.22}{\milli\volt} \textbar{} 16\,Bit \SI{76}{\micro\volt}

\section*{Rechenbeispiele}
\begin{itemize}
\item Teiler $2\times\SI{1}{\mega\ohm}$ an \SI{10}{\volt}: DMM \SI{4.76}{\volt}, Zeiger (\SI{200}{\kilo\ohm}) \SI{1.43}{\volt}
\item \SI{5}{\volt}/\SI{25}{\ohm} + \SI{1}{\ohm} Shunt: \SI{192}{\milli\ampere} statt 200
\item SAR 4\,Bit, \SI{16}{\volt}, \SI{11.3}{\volt} $\to 1011_2 = 11$
\item $\pm$(1\,\% + 2 Digits) @ \SI{100.0}{\volt} $\to$ 98,8\dots\SI{101.2}{\volt}
\item 3½ Stellen = 1999 Counts; 6000 Counts @ \SI{6}{\volt} $\to$ \SI{1}{\milli\volt}
\end{itemize}

\section*{Zeitplan}
\small\begin{tabular}{@{}lll@{}}
Folie & 15 Min & 10 Min\\ \midrule
1--3 Einstieg & 1:45 & 1:15\\
4--8 Messgeräte & 5:15 & 2:45\\
9--13 ADC & 5:30 & 3:40\\
14--15 Abschluss & 2:00 & 1:30\\
\end{tabular}\normalsize
\end{multicols}
\end{document}
